\documentclass[11pt]{article}
\usepackage[margin=1in]{geometry}
\usepackage{graphicx}
\usepackage{amsmath,amssymb}
\usepackage{booktabs}
\usepackage{url}
\usepackage[numbers]{natbib}
\usepackage{caption}
\usepackage{subcaption}
\usepackage{xcolor}
\usepackage[colorlinks=true,linkcolor=blue,citecolor=blue,urlcolor=blue]{hyperref}

\title{Terrain-Conditioned Diffusion for Generating Street-Level
Expansion Layouts of Small Towns}

\author{Eshan Iyer\\
\small Independent researcher\\
\small \texttt{Code: \url{https://github.com/eshankiyer/terrain-urban-diffusion}}}

\date{July 2026}

\begin{document}
\maketitle

\begin{abstract}
Generative models for urban layouts have concentrated on large,
grid-dominated cities, typically in the United States, and almost
universally ignore the terrain the city sits on. Yet for the small towns
that make up most settlements worldwide---and especially for towns in
hilly regions of Europe and Asia---topography is a first-order constraint
on where streets and buildings can go. We present a conditional denoising
diffusion model that generates street-level expansion layouts (road
rasters and built-up area) for small towns, conditioned jointly on
terrain (elevation and slope) and the existing settlement footprint. The
model is trained on windows around approximately sixty towns in hilly
terrain across Europe and Asia, assembled automatically from
OpenStreetMap vectors and open elevation tiles, using a concentric-erosion
proxy to recover core/expansion training pairs from present-day
footprints. A single trained model supports two inference modes: a
\emph{real mode}, which proposes expansions for actual towns, and a
\emph{sandbox mode}, in which the model is applied iteratively to an
arbitrary heightmap so that a settlement grows step by step, in the manner
of a city-building game. On held-out towns, generated expansions track the
slope-occupancy statistics of real growth substantially better than a
slope-blind baseline while maintaining road connectivity to the existing
network. We release code, the town list, and the training pipeline.
\textcolor{red}{[RESULTS-PLACEHOLDER: final numbers inserted from
evaluation run.]}
\end{abstract}

\section{Introduction}

Most generative work on urban form asks the question: what does a city
look like? This paper asks a narrower and, we argue, more practically
useful question: given a particular small town, on particular terrain,
what would its next increment of growth plausibly look like?

The distinction matters for three reasons. First, most of the world's
settlements are not large cities. Small towns dominate by count, and their
growth decisions---where a new residential street goes, which hillside
gets terraced---are made with far less planning capacity than a
metropolitan government has. A tool that proposes terrain-respecting
expansion options is more useful at this scale, not less. Second, terrain
is the binding constraint precisely where planning capacity is lowest:
the hill towns of Nepal, the Italian Apennines, or the Western Ghats
cannot borrow street patterns from Chicago. Third, existing generative
models have a documented geographic bias: the most capable urban layout
generators are trained on large, mostly American cities
\citep{he2025stepwise, xie2023citygen, he2023globalmapper}, whose gridded
morphology transfers poorly to organically grown towns; where non-US
cities have been included, the models have measurably struggled with
non-grid street orientation \citep{jiang2025dcgan}.

Terrain-aware settlement growth modelling is not itself new. The SLEUTH
family of cellular automata has used slope as a growth driver since the
1990s \citep{clarke1998sleuth}, and modern deep-learning land-use-change
models routinely include elevation among their covariates
\citep{chughtai2024casurvey}. But these models predict growth
\emph{probability} at coarse resolution: their output is a per-pixel
likelihood of urbanisation, not a design. They cannot say where the
street goes. Conversely, recent diffusion-based layout generators produce
street-level designs but condition on land use, text, or nothing at all
--- not on topography \citep{he2025stepwise, espinosa2023scotland,
wang2025satellite}. Diffusion models that do engage with terrain generate
the terrain itself \citep{borne2025mesa, hu2023tdn}, not settlements upon
it. The specific combination this paper addresses---street-level
generative layout design, conditioned on terrain, for small towns, trained
on European and Asian data---is, to our knowledge, unoccupied.

Our contributions are:
\begin{enumerate}
\item A fully automatic data pipeline that assembles terrain-conditioned
expansion training pairs for arbitrary towns from OpenStreetMap and open
elevation tiles, requiring no manual annotation and no API keys, together
with a curated list of $\sim$60 training and 8 held-out small towns in
hilly terrain across Europe and Asia.
\item A concentric-erosion proxy that recovers (core, expansion-ring)
supervision pairs from a single present-day snapshot, sidestepping the
scarcity of historical street-level data for small towns; we state its
assumptions and failure modes explicitly.
\item A compact ($\sim$13M parameter) conditional diffusion model that
generates joint road/built-up expansion rasters, and an evaluation
protocol measuring slope compliance, road connectivity, and built-up
contiguity against a slope-blind baseline on held-out towns.
\item A demonstration that one trained model supports both proposal
generation for real towns and iterative, game-like growth of synthetic
settlements on arbitrary heightmaps, which we discuss as the
decision-making layer of a prospective automated city-building agent.
\end{enumerate}

\section{Related Work}

\subsection{Generative urban layout models}
Procedural methods based on shape grammars and production rules have long
been used to synthesise street networks \citep{parish2001procedural,
beirao2010grammar}, and remain the standard in game tooling; their
weakness is the manual rule-crafting required per style. Learning-based
approaches replaced rules with data: Neural Turtle Graphics models road
layouts autoregressively \citep{chu2019ntg}, GlobalMapper generates
building layouts for arbitrary blocks with graph attention
\citep{he2023globalmapper}, and CityGen \citep{xie2023citygen} and
CityDreamer \citep{xie2025citydreamer} produce unbounded city layouts
with diffusion. Closest to our setting, He et al.\ propose a stepwise
multimodal-diffusion urban design framework trained on Chicago and New
York \citep{he2025stepwise}, and explicitly note the restriction to
US-grid morphology as a limitation. GAN-based urban form models trained
on multiple cities report degraded fidelity on non-grid cities
\citep{jiang2025dcgan}. None of these condition on terrain.

\subsection{Urban growth and land-use-change models}
SLEUTH \citep{clarke1998sleuth} takes slope, land use, exclusions, urban
extent, transport and hillshade as inputs to a cellular automaton and
remains widely applied \citep{nigeria2024sleuth}; extensions couple CA
with support vector machines \citep{rienow2015svm} or deep networks
\citep{chughtai2024casurvey}. Deep sequence models (ConvLSTM) forecast
urban growth from satellite time series \citep{boulila2024convlstm}.
These models incorporate terrain but output coarse urbanisation
probabilities rather than street-level layouts; they simulate
\emph{whether} a cell urbanises, not \emph{what is built}.

\subsection{Terrain and geospatial generation}
Diffusion models generate realistic terrain from text or sketches
\citep{borne2025mesa, hu2023tdn, terrafusion2025}, and
map-or-primitive-conditioned diffusion synthesises satellite imagery
\citep{espinosa2023scotland, wei2026terradit}. GPS coordinates have been
used as conditioning for image generation \citep{feng2025gps}. This line
of work treats terrain as the \emph{output} or the imagery as the target;
we invert the relationship, treating terrain as a constraint under which
settlement structure must be generated.

\subsection{Generative models as game agents}
Automating construction in city-building games has mostly meant
procedural map generation. Playing such a game requires perception,
decision-making, and actuation layers. We position our model as the
decision-making layer---current city state and heightmap in, next
construction increment out---and discuss integration via game modding
APIs in Section~\ref{sec:applications}.

\section{Method}

\subsection{Problem formulation}
Let $E \in \mathbb{R}^{H \times W}$ be an elevation raster over a
window of $\sim$1.9\,km ($H{=}W{=}128$ at 15\,m/px), $S$ its slope, and
$(R_c, B_c)$ binary rasters of the \emph{core} road network and built-up
area. The model learns the conditional distribution
$p(R_r, B_r \mid E, S, R_c, B_c)$ over \emph{ring} (expansion) roads and
built-up area. Conditioning channels are concatenated: elevation is
z-scored per window and clipped to $\pm 3\sigma$; slope is normalised by
$30^\circ$. Targets are mapped to $\{-1, 1\}$.

\subsection{Growth supervision by concentric erosion}
\label{sec:erosion}
Historical street-level data for small towns is scarce and inconsistent
across countries. We instead construct supervision from present-day
snapshots: the observed built-up mask $B$ is eroded with a disc of radius
$\rho = 10$\,px ($\approx$150\,m) to produce a pseudo-historical core
$B_c$, and the ring $B_r = B \setminus B_c$ becomes the generation
target; ring roads are defined analogously. The proxy assumes
predominantly outward, concentric growth. This assumption is reasonable
for compact hill towns, whose steep surroundings suppress leapfrog
development, but it is wrong for linear valley towns that grow along a
single axis and for towns whose growth is dominated by infill. We keep
the proxy because it makes the pipeline fully automatic, and we flag the
resulting bias in Section~\ref{sec:limitations}. Nothing in the
architecture depends on it: given true temporal pairs (e.g.\ from
historical cadastral data or multi-epoch built-up products), the same
model trains unchanged.

\subsection{Architecture and training}
A U-Net with three resolution levels (64/128/256 channels), residual
blocks, sinusoidal timestep embeddings, and self-attention at the two
coarsest resolutions predicts the added noise $\epsilon$ under a cosine
noise schedule with $T{=}1000$ steps \citep{ho2020ddpm,
nichol2021improved}. The 4 conditioning channels are concatenated to the
2 noisy target channels at the input ($\sim$13M parameters). We train
with AdamW, cosine learning-rate decay, gradient clipping at 1.0, and an
EMA of the weights (decay 0.999); sampling uses 50 DDIM steps
\citep{song2021ddim}. Each town window is augmented by the dihedral
group (8 orientations) and small centre jitter.

\subsection{Two inference modes}
\textbf{Real mode.} For a real town, the pipeline fetches current
elevation and OSM rasters, erodes to a core exactly as in training, and
samples an expansion conditioned on the true terrain.
\textbf{Sandbox mode.} For an arbitrary heightmap (fractal noise here; a
game-exported heightmap in principle), a seed settlement of a few pixels
is placed on low-slope ground and the model is applied iteratively: the
union of the previous state and the generated ring becomes the next
core. Each round is one ``turn'' of automated play.

\section{Experiments}

\subsection{Data}
$\sim$60 training towns and 8 held-out towns (population roughly
5k--60k) in hilly regions of Europe (Alps, Apennines, Iberia, Balkans)
and Asia (Japan, the Indian Himalaya and Western Ghats, Nepal, Vietnam,
Indonesia, Turkey, the Caucasus), listed in the repository. Windows with
insufficient settlement are dropped automatically.
\textcolor{red}{[RESULTS-PLACEHOLDER: final sample counts.]}

\subsection{Metrics and baseline}
On each held-out town we compare the generated ring to the real ring:
\textbf{Slope KS} --- Kolmogorov--Smirnov distance between the slope
distributions under generated vs.\ real ring built-up pixels (lower is
better; measures whether the model builds on the kinds of slopes the
town actually built on);
\textbf{Road connectivity} --- fraction of generated road pixels
8-connected to the core network;
\textbf{Built contiguity} --- fraction of generated ring pixels adjacent
to at least two settlement pixels;
\textbf{Volume ratio} --- generated/real ring area.
The baseline scatters the same built area uniformly at random over the
window (slope-blind), a deliberately weak control that isolates whether
the model uses terrain at all.

\subsection{Results}
\textcolor{red}{[RESULTS-PLACEHOLDER: table of metrics, model vs
baseline, overall and split by Europe/Asia; loss curve; qualitative
figures for real-mode case studies; sandbox growth sequence figure;
discussion of what the numbers show, written after the numbers exist.]}

\section{Applications: towards automated city-building agents}
\label{sec:applications}
The sandbox mode is a decision-making layer for game automation: at each
turn it maps (heightmap, current city) to a construction increment. A
complete agent for a title such as Cities: Skylines~2 additionally needs
perception (reading the game's heightmap and current road/zone state,
available through its modding API) and actuation (executing placements,
also exposed to mods). Because our model consumes and produces rasters,
the integration surface is thin. We do not implement it here; we note
that raster-to-game-geometry conversion (vectorising road rasters into
placeable splines) is the main engineering step.

\section{Limitations}
\label{sec:limitations}
The concentric-erosion proxy biases training toward ring-like growth
(Section~\ref{sec:erosion}); multi-epoch built-up products such as GHSL
would provide true temporal supervision at coarser spatial detail. OSM
completeness varies across Asia, under-representing buildings in some
towns; our built-up masks partially compensate with land-use polygons.
Rasters at 15\,m/px cannot represent parcel-level structure. The
evaluation baseline is deliberately weak; comparison against SLEUTH-class
CA models and layout generators retrofitted with terrain channels is
future work. Generated layouts are proposals, not plans: no zoning law,
hydrology, hazard, or land-ownership information enters the model.
\textcolor{red}{[RESULTS-PLACEHOLDER: any additional limitations
observed in the actual training run.]}

\section{Conclusion}
We presented a terrain-conditioned diffusion model for street-level
expansion of small towns, trained on European and Asian data assembled
fully automatically, evaluated on held-out towns, and usable both as a
proposal tool for real settlements and as an iterative growth engine on
arbitrary heightmaps.

\section*{AI assistance disclosure}
The code, experiments, and text of this paper were produced with
substantial assistance from an AI system (Anthropic Claude), directed
and reviewed by the author. This disclosure is made in accordance with
venue policies on AI-assisted research.

\bibliographystyle{unsrtnat}
\begin{thebibliography}{99}

\bibitem[Clarke and Gaydos(1998)]{clarke1998sleuth}
K.~C. Clarke and L.~J. Gaydos.
\newblock Loose-coupling a cellular automaton model and GIS: long-term urban
growth prediction for San Francisco and Washington/Baltimore.
\newblock \emph{International Journal of Geographical Information Science},
12(7):699--714, 1998.

\bibitem[He et~al.(2025)]{he2025stepwise}
M.~He, Y.~Liang, S.~Wang, Y.~Zheng, Q.~Wang, D.~Zhuang, L.~Tian, and J.~Zhao.
\newblock Generative AI for urban design: A stepwise approach integrating
human expertise with multimodal diffusion models.
\newblock \emph{arXiv preprint arXiv:2505.24260}, 2025.

\bibitem[Xie et~al.(2023)]{xie2023citygen}
Z.~Xie et~al.
\newblock CityGen: Infinite and controllable city layout generation.
\newblock \emph{arXiv preprint arXiv:2312.01508}, 2023.

\bibitem[He and Aliaga(2023)]{he2023globalmapper}
L.~He and D.~Aliaga.
\newblock GlobalMapper: Arbitrary-shaped urban layout generation.
\newblock In \emph{ICCV}, 2023.

\bibitem[Jiang et~al.(2025)]{jiang2025dcgan}
F.~Jiang et~al.
\newblock A generative urban form design framework based on deep
convolutional GANs and landscape pattern metrics for sustainable renewal in
highly urbanized cities.
\newblock \emph{Sustainability}, 17(10):4548, 2025.

\bibitem[Chughtai et~al.(2024)]{chughtai2024casurvey}
A.~H. Chughtai et~al.
\newblock Integrating cellular automata with deep learning for urban growth
modelling: a review.
\newblock 2024.

\bibitem[Espinosa and Crowley(2023)]{espinosa2023scotland}
M.~Espinosa and E.~J. Crowley.
\newblock Generate your own Scotland: Satellite image generation conditioned
on maps.
\newblock \emph{arXiv preprint arXiv:2308.16648}, 2023.

\bibitem[Wang et~al.(2025)]{wang2025satellite}
Q.~Wang et~al.
\newblock Generative AI for urban planning: Synthesizing satellite imagery
via diffusion models.
\newblock \emph{arXiv preprint arXiv:2505.08833}, 2025.

\bibitem[Borne-Pons et~al.(2025)]{borne2025mesa}
P.~Borne-Pons, M.~Czerkawski, R.~Martin, and R.~Rouffet.
\newblock MESA: Text-driven terrain generation using latent diffusion and
global Copernicus data.
\newblock \emph{arXiv preprint arXiv:2504.07210}, 2025.

\bibitem[Hu et~al.(2023)]{hu2023tdn}
Z.~Hu, K.~Hu, C.~Mo, L.~Pan, and Z.~Wang.
\newblock Terrain diffusion network: Climatic-aware terrain generation with
geological sketch guidance.
\newblock \emph{arXiv preprint arXiv:2308.16725}, 2023.

\bibitem[Parish and M{\"u}ller(2001)]{parish2001procedural}
Y.~I.~H. Parish and P.~M{\"u}ller.
\newblock Procedural modeling of cities.
\newblock In \emph{SIGGRAPH}, 2001.

\bibitem[Beir{\~a}o et~al.(2010)]{beirao2010grammar}
J.~Beir{\~a}o, J.~Duarte, and R.~Stouffs.
\newblock Implementing a generative urban design model: Grammar-based design
patterns for urban design.
\newblock In \emph{eCAADe}, 2010.

\bibitem[Chu et~al.(2019)]{chu2019ntg}
H.~Chu et~al.
\newblock Neural turtle graphics for modeling city road layouts.
\newblock In \emph{ICCV}, 2019.

\bibitem[Xie et~al.(2025)]{xie2025citydreamer}
H.~Xie, Z.~Chen, F.~Hong, and Z.~Liu.
\newblock CityDreamer4D: Compositional generative model of unbounded 4D
cities.
\newblock \emph{arXiv preprint arXiv:2501.08983}, 2025.

\bibitem[Agyemang et~al.(2024)]{nigeria2024sleuth}
F.~S.~K. Agyemang et~al.
\newblock Assessing the suitability of the SLEUTH cellular automata model for
capturing heterogeneous urban growth in developing cities.
\newblock \emph{PLOS ONE}, 2024.

\bibitem[Rienow and Goetzke(2015)]{rienow2015svm}
A.~Rienow and R.~Goetzke.
\newblock Supporting SLEUTH -- enhancing a cellular automaton with support
vector machines for urban growth modeling.
\newblock \emph{Computers, Environment and Urban Systems}, 49:66--81, 2015.

\bibitem[Boulila et~al.(2024)]{boulila2024convlstm}
W.~Boulila et~al.
\newblock A comparative analysis on the use of a cellular automata Markov
chain versus a convolutional LSTM model in forecasting urban growth using
Sentinel 2A images.
\newblock \emph{Journal of Land Use Science}, 2024.

\bibitem[Ho et~al.(2020)]{ho2020ddpm}
J.~Ho, A.~Jain, and P.~Abbeel.
\newblock Denoising diffusion probabilistic models.
\newblock In \emph{NeurIPS}, 2020.

\bibitem[Nichol and Dhariwal(2021)]{nichol2021improved}
A.~Nichol and P.~Dhariwal.
\newblock Improved denoising diffusion probabilistic models.
\newblock In \emph{ICML}, 2021.

\bibitem[Song et~al.(2021)]{song2021ddim}
J.~Song, C.~Meng, and S.~Ermon.
\newblock Denoising diffusion implicit models.
\newblock In \emph{ICLR}, 2021.

\bibitem[Wei et~al.(2026)]{wei2026terradit}
B.~Wei, S.~Sastry, D.~Cher, E.~Xing, and N.~Jacobs.
\newblock TerraDiT-$\Omega$: Unified spatial control for satellite image
synthesis with any geospatial primitive.
\newblock \emph{arXiv preprint arXiv:2606.31029}, 2026.

\bibitem[Feng et~al.(2025)]{feng2025gps}
C.~Feng, Z.~Chen, A.~Holynski, A.~A. Efros, and A.~Owens.
\newblock GPS as a control signal for image generation.
\newblock \emph{arXiv preprint arXiv:2501.12390}, 2025.

\bibitem[TerraFusion(2025)]{terrafusion2025}
K.~Sasaki et~al.
\newblock TerraFusion: Joint generation of terrain geometry and texture
using latent diffusion models.
\newblock \emph{arXiv preprint arXiv:2505.04050}, 2025.

\end{thebibliography}

\end{document}
